% 多圈重整化、算符混合与虚时边界条件的集中推导。

\begin{derivation}[从\eqnrefs{eq:rbar-operation-dp68}到\eqnrefs{eq:r-operation-dp11}]
设两圈振幅为 $I_\Gamma$，其中 $\gamma$ 是一个对数发散的单圈真子图。对子图定义

\begin{equation*}
 C(\gamma)=-K I_\gamma.
\end{equation*}
在只有这一个发散真子图的情形，正文\eqnrefs{eq:rbar-operation-dp68} 化为
\begin{align*}
 \bar R I_\Gamma
 &=I_\Gamma+C(\gamma)I_{\Gamma/\gamma}\\
 &=I_\Gamma-(K I_\gamma)I_{\Gamma/\gamma}.
\end{align*}
第二项在每一个外层动量处抵消子图的大动量局域极点，因此 $\bar R I_\Gamma$ 只可能还保留整张图共同放大时的整体局域发散。正文\eqnrefs{eq:r-operation-dp11} 已把这一剩余整体发散的局域反项写成 $C(\Gamma)=-K\bar R I_\Gamma$。将该反项插回后便有
\begin{align*}
 R I_\Gamma
 &=\bar R I_\Gamma+C(\Gamma)\\
 &=(1-K)\bar R I_\Gamma,
\end{align*}
即正文\eqnrefs{eq:r-operation-dp11}。多个互不重叠或彼此嵌套的真子图只是在\eqnrefs{eq:rbar-operation-dp68} 中增加允许的子图集合，随后同样只对剩余整体发散作最后一次局域减除。
\end{derivation}

\begin{derivation}[从\eqnrefs{eq:composite-operator-mixing-dp11}到\eqnrefs{eq:wilson-rg-dp11}]
由

\begin{equation*}
 \mathcal O_B=Z^{\mathcal O}\mathcal O_R
\end{equation*}
并固定裸理论参数，
\begin{equation*}
 0=\mu\frac{\dd Z^{\mathcal O}}{\dd\mu}\mathcal O_R
 +Z^{\mathcal O}\mu\frac{\dd\mathcal O_R}{\dd\mu}.
\end{equation*}
左乘 $(Z^{\mathcal O})^{-1}$，并定义
\begin{equation*}
 \gamma^{\mathcal O}
 =(Z^{\mathcal O})^{-1}\mu\frac{\dd Z^{\mathcal O}}{\dd\mu},
\end{equation*}
得到
\begin{equation*}
 \mu\frac{\dd\mathcal O_{i,R}}{\dd\mu}
 =-\gamma_{ij}^{\mathcal O}\mathcal O_{j,R}.
\end{equation*}
有效作用量中的同一局域项写成 $C_i\mathcal O_i$。物理量不依赖任意重整化尺度，所以
\begin{align*}
 0
 &=\mu\frac{\dd}{\dd\mu}(C_i\mathcal O_i)\\
 &=\left(\mu\frac{\dd C_i}{\dd\mu}\right)\mathcal O_i
 -C_i\gamma_{ij}^{\mathcal O}\mathcal O_j.
\end{align*}
把第二项的哑指标交换到与第一项相同的算符基，逐个比较 $\mathcal O_i$ 的系数即得正文\eqnrefs{eq:wilson-rg-dp11}。
\end{derivation}

\begin{derivation}[从\eqnrefs{eq:matsubara-boundary-dp11}到\eqnrefs{eq:matsubara-frequencies-dp11}]
玻色 Fourier 模写成 $\ee^{-\ii\omega_n\tau}$。周期条件给

\begin{equation*}
 \ee^{-\ii\omega_n\beta_T\hbar}=1,
\end{equation*}
因此
\begin{equation*}
 \omega_n^{(B)}=\frac{2\pi n}{\beta_T\hbar},
 \qquad n\in\mathbb Z.
\end{equation*}
费米场满足反周期条件，于是
\begin{equation*}
 \ee^{-\ii\omega_n\beta_T\hbar}=-1,
\end{equation*}
从而
\begin{equation*}
 \omega_n^{(F)}=\frac{(2n+1)\pi}{\beta_T\hbar}.
\end{equation*}
两者合起来就是正文\eqnrefs{eq:matsubara-frequencies-dp11}。在区间 $0\le\tau<\beta_T\hbar$ 上，这些模式满足
\begin{equation*}
 \frac1{\beta_T\hbar}\int_0^{\beta_T\hbar}\dd\tau\,
 \ee^{\ii(\omega_n-\omega_m)\tau}=\delta_{nm},
\end{equation*}
因此零温连续 Fourier 测度相应变为正文\eqnrefs{eq:matsubara-loop-replacement-dp11} 的离散求和。
\end{derivation}
